\chapter{User Stories}
\ChapterLabel{UserStories}
\index{user story}
\index{acceptance criterion}
\section{Chapter Responsibility}
\label{Section:UserStoriesResponsibility}
Capture who needs a result and why that result has value.

\noindent\textbf{Expected result:} A reviewable statement of stakeholder value and acceptance intent.

\section{SP800Project Status}
\label{Section:SP800StatusUserStories}
The original stories US-001 through US-003 are inherited template examples.
The project stories \UserStoryRef{US-004} through \UserStoryRef{US-007}
are elaborated drafts based on the instructor feedback supplied by A. Crone.
Each states a capability, outcome, operating context, and acceptance question.
Priorities and detailed requirement definitions remain TBD; these stories
are not approved requirements. A. Crone owns the unresolved decisions;
decision dates are TBD.
See \hyperref[Section:PendingStoryRequirements]{the pending requirement
placeholders} and \hyperref[Section:SP800OpenQuestions]{the project open questions}.

\section{How to Use This Chapter}
Capture stakeholder goals in accessible language. Stories express value and
acceptance intent; precise obligations belong in Requirements. Give each story
a stable ID, a priority, and links to one or more requirements that implement it.

Stakeholder responsibilities, concerns, and expected evidence are maintained in
\hyperref[Table:StakeholderDefinitions]{the stakeholder definitions table}.
Use the expected evidence to refine the draft capability, outcome, context,
and acceptance question for each story.

\section{User Story Register}
\begin{tabularx}{\linewidth}{L{1.9cm}YL{1.8cm}L{3.4cm}}
\caption{User Story Register\label{Table:UserStories}}\\
\toprule
\textbf{ID} & \textbf{Story and Acceptance Criteria} & \textbf{Priority} & \textbf{Requirement(s)}\\
\midrule
\endfirsthead
\multicolumn{4}{l}{\small\itshape User Story Register (continued)}\\
\toprule
\textbf{ID} & \textbf{Story and Acceptance Criteria} & \textbf{Priority} & \textbf{Requirement(s)}\\
\midrule
\endhead
\midrule
\multicolumn{4}{r}{\small\itshape Continued on next page}\\
\endfoot
\bottomrule
\endlastfoot
\UserStoryID{US-001} & As a customer, I want to submit a service request so that the responsible team can act on it. Acceptance: valid data produces a unique request identifier on every supported client. & \gls{must} & \RequirementRef{REQ-F-001}, \RequirementRef{REQ-C-001}\\
\UserStoryID{US-002} & As a customer, I want to view my request status so that I understand its progress. Acceptance: the current status and most recent update appear within the specified response time, and the transition history is retained. & \gls{should} & \RequirementRef{REQ-F-002}, \RequirementRef{REQ-Q-001}, \RequirementRef{REQ-D-001}, \RequirementRef{REQ-C-001}\\
\UserStoryID{US-003} & As an operations analyst, I want to export request history so that I can perform approved analysis. Acceptance: an authorized actor prepares complete history within five minutes in the approved versioned format without changing source records. & \gls{would} & \RequirementRef{REQ-B-001}, \RequirementRef{REQ-F-003}, \RequirementRef{REQ-I-001}, \RequirementRef{REQ-S-001}, \RequirementRef{REQ-D-001}, \RequirementRef{REQ-C-001}\\
\UserStoryID{US-004} & As an \hyperref[Stakeholder:Operator]{operator}, I need to see current sensor values together with their availability, validity, and age so that I can decide whether to use a reading. \par\textit{Draft stakeholder need.} \par\textbf{Operating context:} During an acquisition session, the operator assesses published readings while the sensor may be producing valid data, reporting invalid data, becoming unavailable, or leaving an older reading as the most recent observation. \par\textbf{Acceptance question:} Given representative available, unavailable, valid, invalid, and aging readings, can the operator identify the current sensor value and correctly interpret its availability, validity, and age to decide whether to use it? & TBD & \RequirementRef{REQ-F-004}, \RequirementRef{REQ-F-008}, \RequirementRef{REQ-Q-002}, \RequirementRef{REQ-Q-003}, \RequirementRef{REQ-Q-006}, \RequirementRef{REQ-Q-007}, \RequirementRef{REQ-Q-008}, \RequirementRef{REQ-F-010}, \RequirementRef{REQ-I-004}\\
\UserStoryID{US-005} & As an \hyperref[Stakeholder:AcquisitionDeveloper]{acquisition-developer}, I need sensor acquisition to continue while previously acquired samples are processed and delivered so that processing and delivery do not unnecessarily interrupt collection of new samples. \par\textit{Draft stakeholder need.} \par\textbf{Operating context:} The host application receives successive samples while earlier samples remain in processing or delivery. Workload, buffering capacity, and behavior when processing or delivery cannot keep up remain to be specified. \par\textbf{Acceptance question:} Under an agreed workload, can the acquisition-developer demonstrate that new samples continue to be acquired while prior samples are processed and delivered, with sample ordering and any loss or interruption accounted for? & TBD & \RequirementRef{REQ-F-005}, \RequirementRef{REQ-F-008}, \RequirementRef{REQ-I-002}, \RequirementRef{REQ-Q-002}, \RequirementRef{REQ-Q-004}, \RequirementRef{REQ-Q-005}, \RequirementRef{REQ-F-009}, \RequirementRef{REQ-Q-008}\\
\UserStoryID{US-006} & As a \hyperref[Stakeholder:FieldEngineer]{field-engineer}, I need the sensor-acquisition subsystem to run on macOS, Windows, and Linux so that I can deploy and support it across the operating environments used in the field. \par\textit{Draft stakeholder need.} \par\textbf{Operating context:} The field-engineer installs, configures, and operates the host subsystem on each supported platform with its required acquisition interface and driver. Supported operating-system versions and driver configurations remain to be specified. \par\textbf{Acceptance question:} On each agreed macOS, Windows, and Linux configuration, can the field-engineer follow the documented setup procedure and demonstrate the same agreed acquisition, publication, and status behavior? & TBD & \RequirementRef{REQ-F-006}\\
\UserStoryID{US-007} & As a \hyperref[Stakeholder:MissionIntegrator]{mission-integrator}, I need a documented sensor data and status interface so that a sensor or downstream consumer can be replaced without redesigning the connection. \par\textit{Draft stakeholder need.} \par\textbf{Operating context:} During integration or replacement of a sensor-side implementation or downstream consumer, participants use the documented interface contract. The contract must explain data representation, field meanings, status behavior, timing semantics, and compatibility expectations. OMS applicability remains a separate decision. \par\textbf{Acceptance question:} Given the documented contract and representative data and status examples, can the mission-integrator substitute a conforming sensor-side implementation or consumer without redesigning the connection, and identify an incompatible replacement before use? & TBD & \RequirementRef{REQ-F-007}, \RequirementRef{REQ-I-002}, \RequirementRef{REQ-I-003}, \RequirementRef{REQ-Q-003}, \RequirementRef{REQ-F-009}, \RequirementRef{REQ-Q-007}, \RequirementRef{REQ-F-010}, \RequirementRef{REQ-I-004}\\
\UserStoryID{US-008} & As a \hyperref[Stakeholder:ProductOwner]{product owner}, I need the system and software architecture design to be compatible with OMS so that it can be integrated into future mission systems. \par\textbf{Operating context:} During architecture development and technical review, the product owner evaluates readiness for future OMS integration. \par\textbf{Acceptance question:} Does the project architecture documentation demonstrate conformance to the applicable provisions of the selected OMS Version 2.5 baseline, with traceable evidence and documented gaps? \par\textit{Draft need; authoritative revision, applicability, and assessment details TBD. Owner: A. Crone; decision date: TBD.} & TBD & \RequirementRef{REQ-C-002}, \RequirementRef{REQ-I-004}\\
\end{tabularx}

\section{Reminder: Complete Story Requirements}
\label{Section:StoryRequirementsReminder}
\noindent\textbf{To do --- A. Crone; decision date: TBD.}
Review the draft capabilities, outcomes, operating contexts, and acceptance
questions in \UserStoryRef{US-004}, \UserStoryRef{US-005},
\UserStoryRef{US-006}, and \UserStoryRef{US-007} with the corresponding
stakeholders; collect the planned evidence in
\hyperref[Table:StakeholderDefinitions]{the stakeholder definitions table}
and assign each priority. Agree the workload, buffering and overload behavior,
supported platform and driver configurations, and interface compatibility
criteria before approving dependent requirements. These values remain TBD;
no numerical performance bound or passing evidence is implied.
Review and resolve TBDs in the linked
\hyperref[Section:PendingStoryRequirements]{candidate requirements}, then
complete test, verification, and validation links in
\hyperref[Table:TraceabilityMatrix]{the traceability register}.

\section{Review Rules}
Confirm that every story identifies a beneficiary, outcome, acceptance
criterion, priority, and at least one linked requirement. Maintain reciprocal
story links in the traceability register.

\section{Story Validation Plans}
\UserStoryRef{US-004} is evaluated by \ValidationRef{VAL-004} with \hyperref[PlannedEvidence:VAL-004]{planned stakeholder evidence}.

\UserStoryRef{US-005} is evaluated by \ValidationRef{VAL-005} with \hyperref[PlannedEvidence:VAL-005]{planned stakeholder evidence}.

\UserStoryRef{US-006} is evaluated by \ValidationRef{VAL-006} with \hyperref[PlannedEvidence:VAL-006]{planned stakeholder evidence}.

\UserStoryRef{US-007} is evaluated by \ValidationRef{VAL-007} with \hyperref[PlannedEvidence:VAL-007]{planned stakeholder evidence}.

\UserStoryRef{US-008} is evaluated by \ValidationRef{VAL-008} with \hyperref[PlannedEvidence:VAL-008]{planned stakeholder evidence}.

